\section{Вариационные методы решения краевых задач и определения собственных значений}
\label{sec:variational-bvp-eigen}

\subsection{От дифференциальной к слабой постановке}

Для пояснения общей вариационной схемы рассмотрим стандартную эллиптическую краевую задачу
\[
    -\nabla\!\cdot\!(k\nabla u)=f\quad\text{в }\Omega,
\]
с условиями Дирихле на $\Gamma_D$ и Неймана на $\Gamma_N$:
\[
    u=0\quad\text{на }\Gamma_D,
    \qquad
    k\frac{\partial u}{\partial n}=g\quad\text{на }\Gamma_N.
\]
Умножим уравнение на пробную функцию $v$, равную нулю на $\Gamma_D$, и
проинтегрируем по частям. Получаем слабую постановку: найти $u\in V$, где
\[
    V=\{v\in H^1(\Omega):v|_{\Gamma_D}=0\},
\]
такую, что
\[
    a(u,v)=l(v)\qquad \forall v\in V,
\]
где
\[
    a(u,v)=\int_{\Omega}k\,\nabla u\cdot\nabla v\,d\Omega,
    \qquad
    l(v)=\int_{\Omega}fv\,d\Omega+
         \int_{\Gamma_N}gv\,d\Gamma.
\]
Слабая форма требует меньшей гладкости решения и является исходной для
вариационных и конечно-элементных методов.

\subsection{Энергетический функционал}

Если билинейная форма $a(\cdot,\cdot)$ симметрична и положительно определена,
то слабая задача эквивалентна минимизации функционала
\[
    J(v)=\frac12a(v,v)-l(v),
    \qquad v\in V.
\]
Действительно, условие стационарности $\delta J(u;w)=0$ для всех $w\in V$
даёт
\[
    a(u,w)-l(w)=0.
\]
В механике этот функционал часто имеет смысл полной потенциальной энергии:
равновесное состояние реализует её минимум среди кинематически допустимых
полей.

\subsection{Метод Ритца}

В методе Ритца бесконечномерное пространство $V$ заменяют конечномерным
подпространством
\[
    V_N=\operatorname{span}\{\varphi_1,\ldots,\varphi_N\}\subset V
\]
и ищут приближение
\[
    u_N=\sum_{j=1}^{N}c_j\varphi_j.
\]
Базисные функции заранее удовлетворяют главным, то есть кинематическим,
граничным условиям. Коэффициенты определяют из условия минимума
$J(u_N)$:
\[
    \frac{\partial J}{\partial c_i}=0,
    \qquad i=1,\ldots,N.
\]
В результате возникает система
\[
    \sum_{j=1}^{N}a(\varphi_j,\varphi_i)c_j=l(\varphi_i),
\]
или в матричной форме
\[
    K\mathbf c=\mathbf f,
    \qquad
    K_{ij}=a(\varphi_j,\varphi_i),
    \quad f_i=l(\varphi_i).
\]
При увеличении числа базисных функций пространство допустимых приближений
расширяется и решение уточняется.

\subsection{Метод Галёркина}

Метод Галёркина применим и тогда, когда явный энергетический функционал не
вводится. Подставляют приближение $u_N\in V_N$ в исходное уравнение и требуют,
чтобы его невязка была ортогональна всем пробным функциям того же пространства.
В абстрактной форме:
\[
    a(u_N,v_N)=l(v_N)
    \qquad \forall v_N\in V_N.
\]
Если $v_N=\varphi_i$, снова получается система $K\mathbf c=\mathbf f$.
Для симметричных потенциальных задач методы Ритца и Галёркина приводят к одной
и той же системе уравнений: Ритц исходит из минимума функционала, а Галёркин
--- из ортогональности невязки.

\subsection{Задача на собственные значения}

Во многих краевых задачах требуется найти не только функцию $u$, но и параметр
$\lambda$, при котором существует ненулевое решение. Слабая постановка имеет
вид
\[
    a(u,v)=\lambda\,b(u,v)
    \qquad \forall v\in V,
    \qquad u\neq0,
\]
где $a$ и $b$ --- симметричные билинейные формы, а $b(u,u)>0$ для $u\neq0$.
Например, для свободных колебаний дискретизированной механической системы
\[
    K\mathbf x=\omega^2 M\mathbf x.
\]
Здесь собственное значение $\lambda=\omega^2$ определяет квадрат собственной
частоты, а собственный вектор --- форму колебаний.

\subsection{Отношение Рэлея и вариационный принцип}

Для любой ненулевой функции $v\in V$ вводят отношение Рэлея
\[
    R(v)=\frac{a(v,v)}{b(v,v)}.
\]
Если $u$ --- собственная функция, то $R(u)=\lambda$. Для самосопряжённой
положительной задачи первое собственное значение характеризуется минимумом
\[
    \lambda_1=\min_{v\in V,\,v\neq0}R(v).
\]
Следующие собственные значения можно получать последовательной минимизацией
при условиях ортогональности к ранее найденным собственным функциям в
скалярном произведении $b(\cdot,\cdot)$:
\[
    \lambda_k=
    \min_{\substack{v\neq0,\\
    b(v,u_j)=0,\;j=1,\ldots,k-1}}R(v).
\]
Это и есть основная вариационная идея определения собственных значений.

Для механической системы
\[
    R(\mathbf x)=\frac{\mathbf x^T K\mathbf x}
                       {\mathbf x^T M\mathbf x}.
\]
Стационарные значения этого отношения являются собственными значениями
обобщённой задачи $K\mathbf x=\lambda M\mathbf x$.

\subsection{Метод Рэлея--Ритца для собственных значений}

Выберем конечномерное пространство $V_N$ и положим
\[
    u_N=\sum_{j=1}^{N}c_j\varphi_j.
\]
Подстановка в вариационную задачу приводит к обобщённой матричной задаче
\[
    K\mathbf c=\lambda M\mathbf c,
\]
где
\[
    K_{ij}=a(\varphi_j,\varphi_i),
    \qquad
    M_{ij}=b(\varphi_j,\varphi_i).
\]
Её ненулевое решение существует, когда
\[
    \det(K-\lambda M)=0.
\]
Таким образом, бесконечномерная спектральная задача заменяется конечной
задачей на собственные значения. Тот же результат получается методом
Галёркина, если требовать ортогональности невязки собственным пробным функциям.

\subsection{Сходимость приближений}

Качество решения определяется выбором пространства $V_N$. Если
последовательность пространств становится всё более полной в $V$, а билинейная
форма задачи обладает стандартными свойствами непрерывности и коэрцитивности,
то галёркинские и ритцевские приближения сходятся к точному решению. Для
спектральной задачи при обогащении пространства приближённые собственные
значения и формы сходятся к собственным значениям и функциям исходного
самосопряжённого оператора.

\subsection{Что важно сказать на экзамене}

Нужно показать общую цепочку:
\[
    \text{краевая задача}
    \;\longrightarrow\;
    \text{слабая форма }a(u,v)=l(v)
    \;\longrightarrow\;
    \text{конечномерное пространство }V_N.
\]
После этого различить Ритца (минимизация энергетического функционала) и
Галёркина (ортогональность невязки). Для собственных значений обязательно
записать
\[
    a(u,v)=\lambda b(u,v),
    \qquad
    R(u)=\frac{a(u,u)}{b(u,u)},
\]
и объяснить, что минимум отношения Рэлея даёт первое собственное значение, а
после дискретизации получается задача $K\mathbf c=\lambda M\mathbf c$.

\newpage
